\section{Orthogonal Bases}
\bx{
\en{
    \item 
    \item 
    \item 
    \item 
    \item 
    \item 
    \item 
    Let $V$ be a finite dimensional vector space with a positive definite scalar product. 
    Let $W$ be a subspace. 
    Show that 
    \begin{equation*}
        V = W + W^{\perp} \quad \text{and} \quad W \cap W^{\perp} = \{ O \}. 
    \end{equation*}
    In the terminology of the preceding chapter, this means that $V$ is the direct sum of $W$ and its orthogonal complement. 
    [Use Theorem 2.3.] 
    \begin{solution}
        \begin{proof}
            To the first identity, let $\{ w_1, \dots, w_r \}$ be an orthogonal basis for $W$. 
            Then we can extend it to an orthogonal basis $\{ w_1, \dots, w_r, u_1, \dots, u_s \}$ for $V$ by Theorem 2.1. 
            Let $U$ be the subspace generated by $u_1, \dots, u_s$. 
            Then $U = W^{\perp}$ by Theorem 2.3. 

            Let $v \in V$. There exist numbers $a_i$ ($i = 1, \dots, r$) and $b_i$ ($i = 1, \dots, s$) such that 
            \begin{equation*}
                v = \sum_{i=1}^r a_iw_i + \sum_{j=1}^s b_ju_j. 
            \end{equation*}
            We can regard the first summation as an element $w$ of $W$ and the second summation as an element $u$ of $U$, which also implies the second sum is an element of $W^{\perp}$. 
            Then 
            \begin{equation*}
                v = w + u, 
            \end{equation*}
            where $w$ and $u$ are elements of $W$ and $W^{\perp}$, respectively. 
            Hence we conclude that 
            \begin{equation*}
                V = W + W^{\perp}.                 
            \end{equation*}

            Next, we come to the second identity. 
            Let $v$ be an element of $W \cap W^{\perp}$, which implicates $v$ is simultaneously an element of $W$ and an element of $W^{\perp}$. 
            Then we have 
            \begin{equation*}
                \langle v, v \rangle = 0. 
            \end{equation*}
            Since $V$ is a finite dimensional vector space with a positive definite scalar product, we have $v = O$ if and only if $\langle v, v \rangle = 0$. 
            Hence we conclude that $W \cap W^{\perp} = \{ O \}$, which completes the proof. 
        \end{proof}
    \end{solution}
    \item 
    In Exercise 7, show that $(W^{\perp})^{\perp} = W$. 
    Why is this immediate from Theorem 2.3? 
    \begin{solution}
        \begin{proof}
            We shall prove that $W \subset (W^{\perp})^{\perp}$ and $(W^{\perp})^{\perp} \subset W$, so $W = (W^{\perp})^{\perp}$. 

            First let $w \in W$. 
            Since $w$ is perpendicular to all elements of $W^{\perp}$, we have 
            \begin{equation*}
                \langle w, v \rangle = 0, 
            \end{equation*}
            where $v$ is an arbitrary element of $W^{\perp}$. 
            By the definition of orthogonal complement, we know $w$ is an element of $(W^{\perp})^{\perp}$. 
            Thus we conclude 
            \begin{equation*}
                W \subset (W^{\perp})^{\perp}. 
            \end{equation*}
            
            Conversely, by Theorem 2.3, we have 
            \begin{equation*}
                \dim W + \dim W^{\perp} = \dim V \quad \text{and} \quad \dim W^{\perp} + \dim (W^{\perp})^{\perp} = \dim V. 
            \end{equation*}
            Substracting the two equations above, then 
            \begin{equation*}
                \dim W = \dim (W^{\perp})^{\perp}. 
            \end{equation*}
            By Theorem 5.6 of Chapter III and the part above, we conclude 
            \begin{equation*}
                W = (W^{\perp})^{\perp}. 
            \end{equation*}
        \end{proof}
        It is immediate from Theorem 2.3, because we can also use the theorem to $W^{\perp}$ to get the result we desired. 
    \end{solution}
    \item 
    \ea{
        \item 
        Let $V$ be the space of symmetric $n \times n$ matrices. 
        For $A, B \in V$ define 
        \begin{equation*}
            \langle A, B \rangle = \tr(AB), 
        \end{equation*}
        where $\tr$ is the trace (sum of diagonal elements). 
        Show that this satisfies all the properties of a positive definite scalar product. 
        (You might already have done this as an exercise in a previous section.) 
        \item 
        Let $W$ be the subspace of matrices $A$ such that $\tr(A) = 0$. 
        What is the dimension of the orthogonal complement of $W$, relative to the scalar product in part (a)? 
        Given an explicit basis for this orthogonal complement. 
    }
    \begin{Solution}
        \ea{
            \item 
            \begin{proof}
                We have proved the first property in the Exercise 8 (a) of Chapter VI before, so we start from the second property. 
                Let $C \in V$. 
                Then by the distributivity of matrices, we have 
                \begin{equation*}
                    \langle A, B + C \rangle = \tr(A(B + C)) = \tr(AB) + \tr(AC) = \langle A, B \rangle + \langle A, C \rangle, 
                \end{equation*}
                which proves the second property. 

                We continue to show the third property. 
                If $x$ is a number, then 
                \begin{equation*}
                    \langle xA, B \rangle = \tr((xA)B) = x \tr(AB) = x \langle A, B \rangle. 
                \end{equation*}
                On the other hand, we have 
                \begin{equation*}
                    \langle A, xB \rangle = \tr(A(xB)) = x \tr(AB) = x \langle A, B \rangle, 
                \end{equation*}
                proving the third property. 

                We finally come to the last property. 
                By the definition of trace, we have 
                \begin{equation*}
                    \langle A, A \rangle = \tr(AA) = \sum_{i=1}^n \sum_{k=1}^n a_{ik}a_{ki} = \sum_{i=1}^n \sum_{k=1}^n a_{ik}^2 \geq 0, 
                \end{equation*}
                because $A$ is symmetric. 
                Only when $a_{ik} = 0$ for all $i, k$ range in $1 \leq i, k \leq n$, i.e. $A = O$, we have $\langle A, A \rangle = 0$. 
                Hence we conclude the whole proof. 
            \end{proof}
            \item 
            Use the dimension formulas. 
            The trace is a linear map, from $V$ to $\mathbf{R}$. 
            Since 
            \begin{equation*}
                \dim V = \dim \Ker \tr + \dim \Im \tr, 
            \end{equation*}
            it follows that $\dim W = \dim V - 1$ so $\dim W^{\perp} = 1$ by Theorem 2.3. 

            Let $I$ be the unit $n \times n$ matrix. 
            Then $\tr(A) = \tr(AI)$ for all $A \in V$, so $\tr(AI) = 0$ for $A \in W$. 
            Hence $I \in W^{\perp}$. 
            Since $W^{\perp}$ has dimension $1$, it follows that $I$ is a basis of $W^{\perp}$. 
            (Simple?) 
        }
    \end{Solution}
    \item 
    Let $A$ be a symmetric $n \times n$ matrix. 
    Let $X, Y \in \mathbf{R}^n$ be eigenvectors for $A$, that is suppose that there exist numbers $a, b$ such that $AX = aX$ and $AY = bY$. 
    Assume that $a \neq b$. 
    Prove that $X, Y$ are perpendicular. 
    \begin{Solution}
        \begin{proof}
            We have $\langle X, AY \rangle = \langle X, bY \rangle = b \langle X, Y \rangle$ and also 
            \begin{equation*}
                \langle X, AY \rangle = \langle {^{t\!}A}X, Y \rangle = \langle AX, Y \rangle = \langle aX, Y \rangle = a\langle X, Y \rangle, 
            \end{equation*}
            because $A$ is symmetric. 
            So $a\langle X, Y \rangle = b\langle X, Y \rangle$. 
            Since $a \neq b$ it follows that $\langle X, Y \rangle = 0$. 
        \end{proof}
    \end{Solution}
}}